\section{Bethlehem, Swaddling Cloths, and the Manger}

\subsection{What the names and narratives establish}

In Biblical Hebrew, \emph{leḥem} means bread and can denote food; the ordinary
synchronic gloss of \emph{Beth Lehem} is therefore ``house of bread.'' Jerome
and Ambrose receive the name in that sense. The lexical gloss does not settle
the historical origin of the place-name, nor does it turn every later bread
reference into authorial wordplay. Matthew and Luke locate Jesus's birth in
Bethlehem to identify the Davidic city and the scriptural horizon of Micah 5.
John 6 never names the city.

Luke twice emphasizes the sign: the child is wrapped in swaddling cloths and
lying in a manger (Lk 2:7, 12), and the shepherds find him so (2:16). The
manger is a feeding place; Luke's noun does not itself say ``altar'' or
``Eucharist.'' The cloths present an actual newborn received into material
poverty and care. A Eucharistic reading can proceed canonically from Luke to
John, but it should be named as reception or synthesis, not smuggled into
Luke's vocabulary.

\subsection{A direct ancient conjunction: Jerome's Paula}

Jerome's \emph{Letter 108} supplies the clearest checked conjunction. Recalling
Paula's pilgrimage to the birth site, Jerome reports her salutation to
Bethlehem as the ``house of bread'' in which was born the Bread who came down
from heaven, and the translation explicitly cross-references John 6:51. The
same scene has Paula contemplate the infant wrapped in cloths and crying in the
manger.

This is real late-fourth-century Latin reception, but its exact attribution
must be preserved. The rhetoric stands in Jerome's memorial letter; within it,
Jerome presents the words as Paula's. It is not a saying of Luke, John, or a
first-century community. Nor does it make the swaddling cloths a Eucharistic
species. It joins sacred geography, Incarnation, and the Johannine Bread from
heaven.

Ambrose makes a related play in Letter 70: Bethlehem, once construed as a
house of wrath in his argument, becomes the house of bread because it received
the Bread that came down from heaven. The typology is therefore not an isolated
modern slogan. Its controlled late-fourth-century form centers on the place
and the incarnate Bread.

\subsection{The manger image and a necessary attribution warning}

A homily printed among the ``Four Homilies'' attributed to Gregory
Thaumaturgus says that the heavenly Bread was laid on the board from which
cattle eat so that human beings living irrationally might receive spiritual
food. It also dwells on the one wrapped in cloths who clothes creation.
This is almost exactly the synthesis commonly repeated today.

It cannot honestly be cited as secure third-century testimony. The patristic
collection that transmits the English text places the four homilies among
``dubious or spurious writings,'' and older editorial discussion already noted
objections to Gregory's authorship. Until its language, manuscript tradition,
and modern critical attribution are controlled, it is evidence for a
pseudonymous or later homiletic tradition, not for Gregory Thaumaturgus.

Chrysostom supplies a safer ancient bridge: the body lying in the manger and
adored by the Magi is the same body received in communion. He does not depend
on the place-name and does not call the manger an altar. His continuity is
Christological---the Eucharistic Body is not another body than the incarnate
one.

